%% kapa: exam=1 round=30 year=2019 period=2 form=A subject=appraisal_law subject_ko=감정평가관계법규 seq=23 number=23
%% source: Q-Net 감정평가사 기출문제 (cst00309) artlSeq=5205205 file=extracted/1008/5205205/2019년 제30회 감정평가사 1차시험 문제지/2019년 제30회 감정평가사 1차시험 문제지/2019년도 제30회 감정평가사 1차 2교시 A형.hwp
%% conversion: kordoc 4.24.51 -> tools/gen1.py (auto); review=proofread-sonnet; flags=none
\begin{question}{23}
국유재산법령상 일반재산에 관한 설명으로 옳지 않은 것은?
\begin{choices}
  \item 일반재산은 대부 또는 처분할 수 있다.
  \item 총괄청은 3년 이상 활용되지 아니한 일반재산을 민간사업자와 공동으로 개발할 수 있다.
  \item 정부출자기업체의 주주 등 출자자에게 해당 기업체의 지분증권을 매각하는 경우에는 일반재산을 수의계약으로 처분할 수 있다.
  \item 정부는 정부출자기업체의 운영체제와 경영구조의 개편을 위하여 필요한 경우에는 일반재산을 현물출자할 수 있다.
  \item 일반재산인 토지와 사유재산인 토지를 교환할 때 쌍방의 가격이 같지 아니하면 그 차액을 금전으로 대납하여야 한다.
\end{choices}
\end{question}
\answer{2}
